% Options for packages loaded elsewhere
\PassOptionsToPackage{unicode}{hyperref}
\PassOptionsToPackage{hyphens}{url}
\documentclass[
]{article}
\usepackage{xcolor}
\usepackage{amsmath,amssymb}
\setcounter{secnumdepth}{-\maxdimen} % remove section numbering
\usepackage{iftex}
\ifPDFTeX
  \usepackage[T1]{fontenc}
  \usepackage[utf8]{inputenc}
  \usepackage{textcomp} % provide euro and other symbols
\else % if luatex or xetex
  \usepackage{unicode-math} % this also loads fontspec
  \defaultfontfeatures{Scale=MatchLowercase}
  \defaultfontfeatures[\rmfamily]{Ligatures=TeX,Scale=1}
\fi
\usepackage{lmodern}
\ifPDFTeX\else
  % xetex/luatex font selection
\fi
% Use upquote if available, for straight quotes in verbatim environments
\IfFileExists{upquote.sty}{\usepackage{upquote}}{}
\IfFileExists{microtype.sty}{% use microtype if available
  \usepackage[]{microtype}
  \UseMicrotypeSet[protrusion]{basicmath} % disable protrusion for tt fonts
}{}
\makeatletter
\@ifundefined{KOMAClassName}{% if non-KOMA class
  \IfFileExists{parskip.sty}{%
    \usepackage{parskip}
  }{% else
    \setlength{\parindent}{0pt}
    \setlength{\parskip}{6pt plus 2pt minus 1pt}}
}{% if KOMA class
  \KOMAoptions{parskip=half}}
\makeatother
\usepackage{longtable,booktabs,array}
\usepackage{calc} % for calculating minipage widths
% Correct order of tables after \paragraph or \subparagraph
\usepackage{etoolbox}
\makeatletter
\patchcmd\longtable{\par}{\if@noskipsec\mbox{}\fi\par}{}{}
\makeatother
% Allow footnotes in longtable head/foot
\IfFileExists{footnotehyper.sty}{\usepackage{footnotehyper}}{\usepackage{footnote}}
\makesavenoteenv{longtable}
\setlength{\emergencystretch}{3em} % prevent overfull lines
\providecommand{\tightlist}{%
  \setlength{\itemsep}{0pt}\setlength{\parskip}{0pt}}
\usepackage{bookmark}
\IfFileExists{xurl.sty}{\usepackage{xurl}}{} % add URL line breaks if available
\urlstyle{same}
\hypersetup{
  hidelinks,
  pdfcreator={LaTeX via pandoc}}

\author{}
\date{}

\begin{document}

\section{Exact Bessel-Function Series for
E(T)}\label{exact-bessel-function-series-for-et}

\emph{Atkinson (1949) · Titchmarsh (1938) · Voronoï (1904) --- Assembled
March 31, 2026}

\begin{center}\rule{0.5\linewidth}{0.5pt}\end{center}

\subsection{1. Notation and Conventions}\label{notation-and-conventions}

\begin{longtable}[]{@{}
  >{\raggedright\arraybackslash}p{(\linewidth - 2\tabcolsep) * \real{0.4211}}
  >{\raggedright\arraybackslash}p{(\linewidth - 2\tabcolsep) * \real{0.5789}}@{}}
\toprule\noalign{}
\begin{minipage}[b]{\linewidth}\raggedright
Symbol
\end{minipage} & \begin{minipage}[b]{\linewidth}\raggedright
Definition
\end{minipage} \\
\midrule\noalign{}
\endhead
\bottomrule\noalign{}
\endlastfoot
(\zeta(s)) & Riemann zeta-function, (s \in \mathbb{C}
\setminus \{1\}) \\
(d(n)) & number of positive divisors of (n \in \mathbb{N}) \\
(\gamma) & Euler--Mascheroni constant \\
(K\_0, K\_1) & modified Bessel functions of the second kind, orders 0,
1; real, positive, (K\_\nu(z) \sim \sqrt{\pi/(2z)},e\^{}\{-z\}) as (z
\to +\infty) \\
(Y\_0, Y\_1) & Bessel functions of the second kind (Neumann), orders 0,
1; (Y\_\nu(z) \sim \sqrt{2/(\pi z)}\sin(z - \nu\pi/2 - \pi/4)) as (z
\to +\infty) \\
(J\_0, J\_1) & Bessel functions of the first kind, orders 0, 1 \\
(H\_0\^{}\{(1)\}) & Hankel function (H\_0\^{}\{(1)\}(z) = J\_0(z) +
iY\_0(z)) \\
(\Delta(x)) & (\sideset{}{'}\sum\_\{n \le x\} d(n) - x(\log x +
2\gamma - 1) - \tfrac{1}{4}) \\
(\Delta\^{}*(x)) & (-\Delta(x) + 2\Delta(2x) - \tfrac{1}{2}\Delta(4x))
(Jutila's variant) \\
(N) & positive integer with (T/(4\pi) \le N \le T/\pi), chosen freely \\
(\sideset{}{'}\sum) & sum in which the last term is halved if (x
\in \mathbb{Z}) \\
\end{longtable}

\textbf{Definition 1 (Mean-square error term).} The function (E(T)) is
defined for (T \ge 2) by

{[} \int\_0\^{}T
\left\textbar{}\zeta!\left(\tfrac{1}{2}+it\right)\right\textbar\^{}2 dt
= T\log\frac{T}{2\pi} + (2\gamma - 1)T + E(T). \tag{*} {]}

Known bounds: (E(T) = O(T\^{}\{1/3+\varepsilon\})), (E(T) =
\Omega(T\^{}\{1/4\})). The exact representation below is independent of
these bounds.

\begin{center}\rule{0.5\linewidth}{0.5pt}\end{center}

\subsection{2. Bessel Function
Prerequisites}\label{bessel-function-prerequisites}

\textbf{Lemma 2.1} (Integral representation of (K\_0)). For (x
\textgreater{} 0):

{[} K\_0(2\pi x) =
\int\_0\^{}\infty \frac{\cos(2\pi x t)}{\sqrt{t^2+1}},dt. \tag{2.1} {]}

More generally, for (\mathrm{Re}(\nu) \textgreater{} -\tfrac{1}{2}) and
(x \textgreater{} 0):

{[} K\_\nu(2\pi x) =
\frac{\sqrt{\pi}\,(2\pi x)^\nu}{2^\nu \Gamma(\nu + \tfrac{1}{2})}
\int\_0\^{}\infty \frac{\cos(2\pi x t)}{(t^2+1)^{\nu+1/2}},dt. \tag{2.2}
{]}

\emph{Proof.} Entry 10.32.9 of the DLMF; substitute (t = \sinh u).

\textbf{Lemma 2.2} (Integral representation of (Y\_0)). For (z
\textgreater{} 0):

{[} Y\_0(z) =
-\frac{2}{\pi}\int\_1\^{}\infty \frac{\cos(zt)}{\sqrt{t^2-1}},dt. {]}

The combination arising in Voronoï's formula is, for (n \ge 1), (x
\textgreater{} 0):

{[} 4K\_0(4\pi\sqrt{nx}) - 2\pi Y\_0(4\pi\sqrt{nx}), {]}

arising from the Mellin--Barnes integral (3.3) below.

\textbf{Lemma 2.3} (Kernel (H(u,n))). For (n \ge 1) and (\mathrm{Re}(u)
\textless{} 2/3), define

{[} H(u,n) :=
2\int\_0\textsuperscript{\infty y}\{-u\}(1+y)\^{}\{u-1\}\cos(2\pi ny),dy.
\tag{2.4} {]}

The integral converges absolutely for (\mathrm{Re}(u) \textless{}
\tfrac{1}{2}) and is analytically continued to (\mathrm{Re}(u)
\textless{} 2/3) by two integrations by parts. For (u =
\tfrac{1}{2}+it):

{[} H!\left(\tfrac{1}{2}+it, n\right) =
2\int\_0\textsuperscript{\infty y}\{-1/2-it\}(1+y)\^{}\{-1/2+it\}\cos(2\pi ny),dy,
\tag{2.5} {]}

absolutely convergent for each (n \ge 1).

\textbf{Lemma 2.4} (Asymptotics). For (z \to +\infty):

{[} K\_\nu(z)
\sim \sqrt{\frac{\pi}{2z}},e\^{}\{-z\}\sum\emph{\{k=0\}\textsuperscript{\infty a\_k(\nu)z}\{-k\},
\qquad Y}\nu(z) \sim \sqrt{\frac{2}{\pi z}}\sin!\left(z -
\frac{\nu\pi}{2} - \frac{\pi}{4}\right) + O(z\^{}\{-3/2\}). {]}

In particular (K\_\nu(4\pi\sqrt{nx})) decays like
(e\^{}\{-4\pi\sqrt{nx}\}) as (n \to \infty) for fixed (x \textgreater{}
0), and (\textbar Y\_\nu(4\pi\sqrt{nx})\textbar{}
\le C(nx)\^{}\{-1/4\}).

\begin{center}\rule{0.5\linewidth}{0.5pt}\end{center}

\subsection{3. Voronoï--Jutila Exact Summation
Formula}\label{voronouxefjutila-exact-summation-formula}

\textbf{Theorem 3.1} (Voronoï 1904; Jutila 1984 --- exact identity). Let
(a \textless{} b) be positive reals and let (f : {[}a,b{]}
\to \mathbb{C}) be smooth. Then

{[} \sideset{}{'}\sum\_\{a \le n \le b\} d(n)f(n) = \int\emph{a\^{}b
f(u)(\log u + 2\gamma),du + \sum}\{n=1\}\^{}\infty d(n)\int\_a\^{}b
f(u)\left[4K_0(4\pi\sqrt{nu}) - 2\pi Y_0(4\pi\sqrt{nu})\right]du.
\tag{3.1} {]}

The series converges absolutely in the (K\_0) terms and conditionally in
the (Y\_0) terms, uniformly on compact subsets of ((0,\infty)). Equation
(3.1) is an exact identity.

\emph{Proof.}

\textbf{Step 1.} By Perron's formula,

{[} \sideset{}{'}\sum\emph{\{n \le x\} d(n) =
\frac{1}{2\pi i}\int}\{c-i\infty\}\^{}\{c+i\infty\}
\frac{\zeta(w)^2\, x^w}{w},dw, \qquad c \textgreater{} 1. {]}

\textbf{Step 2.} Shift the contour to (\mathrm{Re}(w) = \tfrac{1}{2}),
picking up the double pole at (w = 1), which contributes (x\log x +
(2\gamma-1)x), and the simple pole at (w = 0), contributing
(-\tfrac{1}{4}). Apply the functional equation (\zeta(w) =
\chi(w)\zeta(1-w)) on (\mathrm{Re}(w) = \tfrac{1}{2}).

\textbf{Step 3.} The key Mellin--Barnes identity is

{[} \frac{1}{2\pi i}\int\_\{1/2 - i\infty\}\^{}\{1/2+i\infty\}
\left(\frac{\pi n}{u}\right)\textsuperscript{\{-s\}\Gamma(s)}2\cos(\pi s),ds
= 4K\_0(4\pi\sqrt{nu}) - 2\pi Y\_0(4\pi\sqrt{nu}). \tag{3.3} {]}

\textbf{Step 4.} The (K\_0) series converges absolutely (exponential
decay, Lemma 2.4). The (Y\_0) series converges conditionally after Abel
summation using (\Delta(X) = O(X\^{}\{1/3+\varepsilon\})). \(\square\)

\textbf{Corollary 3.3.} For any smooth (f) and any (N \asymp T),
applying (3.1) with (a = 1), (b = N+\tfrac{1}{2}):

{[} \sum\_\{n=1\}\^{}N d(n)f(n) = \int\emph{1\^{}\{N+1/2\} f(u)(\log u +
2\gamma),du + \sum}\{n=1\}\^{}\infty d(n)\int\_1\^{}\{N+1/2\}
f(u)\left[4K_0(4\pi\sqrt{nu}) - 2\pi Y_0(4\pi\sqrt{nu})\right]du.
\tag{3.4} {]}

\begin{center}\rule{0.5\linewidth}{0.5pt}\end{center}

\subsection{\texorpdfstring{4. Titchmarsh's Exact Identity for
(\zeta\^{}2(s))}{4. Titchmarsh's Exact Identity for (\^{}2(s))}}\label{titchmarshs-exact-identity-for-2s}

\textbf{Theorem 4.1} (Titchmarsh 1938 --- exact identity). For (s
\in \mathbb{C}) with (\mathrm{Re}(s) \textgreater{} 0), (s \ne 1), and
any (x \textgreater{} 0):

{[} \zeta(s)\^{}2 = \sum\emph{\{n \le x\}\frac{d(n)}{n^s} -
x\textsuperscript{\{-s\}\sum\emph{\{n \le x\}d(n) +
\frac{(2s - s^2)}{(s-1)^2},x\^{}\{1-s\} +
\frac{s}{s-1},x\^{}\{1-s\}(2\gamma + \log x) + \frac{1}{4}x\^{}\{-s\}
{]} {[} - (4\pi\textsuperscript{2)}s s
\sum}\{n=1\}}\infty \frac{d(n)}{n^{1-s}}\int}\{4\pi\sqrt{nx}\}\textsuperscript{\infty \left[K_1(v) + \frac{\pi}{2}Y_1(v)\right]v}\{-2s\},dv.
\tag{4.1} {]}

This is an exact identity for every (x \textgreater{} 0) and every (s)
with (\mathrm{Re}(s) \textgreater{} 0), (s \ne 1). No approximation is
involved.

\emph{Proof.}

\textbf{Step 1.} Apply the truncated Perron formula to extract the
partial sum (\sum\emph{\{n \le x\}), obtaining an exact contour integral
for the tail (\sum}\{n \textgreater{} x\}).

\textbf{Step 2.} Shift the tail integral to (\mathrm{Re}(w) = 1 -
\mathrm{Re}(s)), picking up the double pole of (\zeta(w)\^{}2/(w-s)) at
(w = 1) and the simple pole at (w = s). Apply the functional equation
(\zeta(w)\^{}2 = \chi(w)\^{}2 \zeta(1-w)\^{}2).

\textbf{Step 3.} The identity

{[} \frac{1}{2\pi i}\int\emph{\{c-i\infty\}\^{}\{c+i\infty\}
\chi(w)\^{}2 (nx)\^{}\{w-s\},dw = (4\pi\textsuperscript{2)}s s,
n\^{}\{s-1\}\int}\{4\pi\sqrt{nx}\}\textsuperscript{\infty \left[K_1(v) + \frac{\pi}{2}Y_1(v)\right]v}\{-2s\},dv
{]}

converts the contour remainder into the Bessel integral in (4.1), via
the Sommerfeld contour integral for (H\_1\^{}\{(1)\}) and one
integration by parts.

\textbf{Step 4.} The (K\_1) term decays as (e\^{}\{-4\pi\sqrt{nx}\}),
giving absolute convergence in (n). The (Y\_1) term converges
conditionally after partial summation on (d(n)). \(\square\)

\textbf{Corollary 4.2.} Setting (s = \tfrac{1}{2}+it) and integrating
over ({[}0,T{]}) converts (4.1) into Atkinson's four-term exact formula
of Theorem 5.1 without any saddle-point approximation.

\begin{center}\rule{0.5\linewidth}{0.5pt}\end{center}

\subsection{5. Atkinson's Exact Pre-Saddle-Point
Representation}\label{atkinsons-exact-pre-saddle-point-representation}

For (y \textgreater{} 0) and (T \ge 2) define:

{[} \Phi(T,y) :=
\frac{\sin\!\left(T\log\frac{1+y}{y}\right)}{\log\frac{1+y}{y}},
\qquad \psi(y) := \frac{1}{\sqrt{y(1+y)}}, \tag{5.1–5.2} {]}

{[} \Psi(T,x) := \int\_0\^{}\infty \Phi(T,y),\psi(y)\cos(2\pi xy),dy.
\tag{5.3} {]}

The exact Stirling remainder is

{[} R\_0(T) := -\frac{T}{2}\log!\left(1 + \frac{1}{4T^2}\right) -
\frac{1}{12T} - \frac{1}{90T^3}, \qquad \textbar R\_0(T)\textbar{}
\le \frac{1}{8T} + \frac{1}{12T} + \frac{1}{90T^3}. \tag{5.4} {]}

For any integer (N) with (T/(4\pi) \le N \le T/\pi):

{[} E\_1(T) := 4\sum\_\{n=1\}\^{}N d(n),\Psi(T,n), \tag{5.5} {]}

{[} E\_2(T) := 4,\Delta\^{}*!\left(N+\tfrac{1}{2}\right)\Psi!\left(T,
N+\tfrac{1}{2}\right), \tag{5.6} {]}

{[} E\_3(T) :=
-\frac{2}{\pi}(\log(N+\tfrac{1}{2})+2\gamma)\int\_0\^{}\infty \frac{\Phi(T,y)\sin(2\pi(N+\tfrac{1}{2})y)}{y\sqrt{y(1+y)}},dy
+
\frac{1}{\pi i}\int\emph{0\^{}\infty \frac{\sin(2\pi(N+\tfrac{1}{2})y)}{y}\int}\{1/2-iT\}\textsuperscript{\{1/2+iT\}\frac{1}{u}!\left(\frac{1+y}{y}\right)}u
du,dy, \tag{5.7} {]}

{[} E\_4(T) :=
4\int\_\{N+1/2\}\^{}\infty \frac{\Delta^*(x)}{x}\int\_0\^{}\infty \frac{\cos(2\pi xy)}{(1+y)\sqrt{y(1+y)}\,\log\frac{1+y}{y}}\left[T\cos\!\left(T\log\frac{1+y}{y}\right) - \left(\tfrac{1}{2} + \frac{1}{\log\frac{1+y}{y}}\right)\sin\!\left(T\log\frac{1+y}{y}\right)\right]dy,dx.
\tag{5.8} {]}

\textbf{Theorem 5.1} (Atkinson's exact formula). For any (T \ge 2) and
any integer (N) with (T/(4\pi) \le N \le T/\pi):

{[} E(T) = E\_1(T) - E\_2(T) + E\_3(T) - E\_4(T) + R\_0(T). \tag{5.9}
{]}

This is an exact identity for every (T \ge 2).

\emph{Proof.}

\textbf{Step 1.} Write (\int\emph{0\^{}T
\textbar{}\zeta(\tfrac{1}{2}+it)\textbar\^{}2,dt = \sum}\{m,n
\ge 1\}(mn)\^{}\{-1/2\}\int\emph{0\^{}T (n/m)\^{}\{it\},dt). The
diagonal (m = n) sums to (T\sum}\{n\} n\^{}\{-1\}), regularised against
the asymptotic (T\log T + (2\gamma-1-\log 2\pi)T) via the Stirling
series for (\log\Gamma(\tfrac{1}{4}+\tfrac{it}{2})); this produces
(R\_0(T)) exactly, with the bound in (5.4).

\textbf{Step 2.} The off-diagonal terms ((m \ne n)) are rearranged by
the substitution (n/m = (1+y)/y) ((y = m/(n-m))). Via Theorem 3.1, the
resulting sum is converted exactly into (\int\emph{0\^{}T
\sum}\{n=1\}\^{}\infty d(n)H(\tfrac{1}{2}+it, n),dt).

\textbf{Step 3.} Exchange (\int\_0\^{}T dt) and (\int\_0\^{}\infty dy)
by Fubini (justified since (y\textsuperscript{\{-1/2\}(1+y)}\{-1/2\}
\in L\^{}1)). The inner (t)-integral evaluates exactly:

{[} \mathrm{Re}\int\_0\^{}T e\^{}\{it\log\frac{1+y}{y}\},dt =
\frac{\sin\!\left(T\log\frac{1+y}{y}\right)}{\log\frac{1+y}{y}} =
\Phi(T,y). {]}

This gives (E\_1(T)) from the truncated sum (\sum\_\{n=1\}\^{}N).

\textbf{Step 4.} The boundary at (n = N+\tfrac{1}{2}) produces (E\_2(T))
exactly.

\textbf{Step 5.} For (n \textgreater{} N): one integration by parts
against (\Delta\^{}*(x)) gives (E\_4(T)) from the integral and (E\_3(T))
from the lower boundary. Both steps are exact. \(\square\)

\textbf{Remark.} The choice of (N \in [T/(4\pi), T/\pi]) is free; (5.9)
holds for every admissible (N). Taking (N \approx T/(2\pi)) minimises
the term count for numerical computation.

\begin{center}\rule{0.5\linewidth}{0.5pt}\end{center}

\subsection{6. Main Theorem: Exact Bessel-Function Series for
(E(T))}\label{main-theorem-exact-bessel-function-series-for-et}

\textbf{Theorem 6.1} (Exact Bessel-function representation). For (T
\ge 2) and any integer (N) with (T/(4\pi) \le N \le T/\pi):

{[} E(T) = 4\sum\_\{n=1\}\^{}N d(n),\Psi\_B(T,n) -
4,\Delta\^{}*!\left(N+\tfrac{1}{2}\right)\Psi\emph{B!\left(T,N+\tfrac{1}{2}\right)
+ E\_3(T) -
4\int}\{N+1/2\}\^{}\infty \frac{\Delta^*(x)}{x},\Xi\_B(T,x),dx +
R\_0(T), \tag{6.1} {]}

where

{[} \Psi\_B(T,n) :=
\int\_0\^{}\infty \frac{\sin\!\left(T\log\frac{1+y}{y}\right)}{\sqrt{y(1+y)}\,\log\frac{1+y}{y}}\cos(2\pi ny),dy,
\tag{6.2} {]}

{[} \Xi\_B(T,x) :=
\int\_0\^{}\infty \frac{\cos(2\pi xy)}{(1+y)\sqrt{y(1+y)}\,\log\frac{1+y}{y}}\left[T\cos\!\left(T\log\frac{1+y}{y}\right) - \left(\tfrac{1}{2} + \frac{1}{\log\frac{1+y}{y}}\right)\sin\!\left(T\log\frac{1+y}{y}\right)\right]dy,
\tag{6.3} {]}

and (R\_0(T)) is given exactly by (5.4). Furthermore, each
(\Psi\_B(T,n)) has the following exact Bessel representations. Under the
substitution (y = \sinh\^{}2(\theta/2)):

{[} \Psi\_B(T,n) =
\int\_0\^{}\infty \frac{\sin(2T\log\coth(\theta/2))}{2\log\coth(\theta/2)}\cos(\pi n\sinh\theta),d\theta,
\tag{6.5} {]}

{[} \Psi\_B(T,n) =
\int\_0\^{}\infty Y\_0(2\pi n\cosh\theta)\sin(T\theta),d\theta,
\tag{6.7a} {]}

where (6.7a) follows from the Mehler--Sonine integral (Y\_0(z) =
\tfrac{2}{\pi}\int\_0\^{}\infty \sin(z\cosh\theta),d\theta) and the
identification (\mathrm{Im}{[}H\_0\^{}\{(1)\}(2\pi n\cosh\theta){]} =
Y\_0(2\pi n\cosh\theta)). The (K\_0) component enters from the contour
rotation (\theta \to \theta + i\pi/2) via (K\_0(z) =
\tfrac{\pi i}{2}H\_0\^{}\{(1)\}(iz)), contributing a correction bounded
by (Ce\^{}\{-2\pi n\}) for each (n).

All equalities in (6.1)--(6.7a) are exact identities. No saddle-point or
stationary-phase step has been applied.

\emph{Proof.} Equations (6.1)--(6.4) restate (5.9) with the
substitutions (E\_1 = 4\sum d(n)\Psi\_B), (E\_2 =
4\Delta\^{}*(N+\tfrac{1}{2})\Psi\_B(T,N+\tfrac{1}{2})), (E\_4 =
4\int(\Delta\^{}*(x)/x)\Xi\_B,dx). For (6.5): substitute (y =
\sinh\^{}2(\theta/2)) in (6.2), computing (\sqrt{y(1+y)} =
\sinh\theta/2), (\log((1+y)/y) = 2\log\coth(\theta/2)), (\cos(2\pi ny) =
\cos(\pi n(\cosh\theta - 1))), (dy = \sinh\theta,d\theta/2). For (6.7a):
use the Mehler--Sonine integral and
(\mathrm{Im}{[}H\_0\^{}\{(1)\}(\cdot){]} = Y\_0(\cdot)). \(\square\)

\textbf{Corollary 6.2} (Double Bessel series for (E\_4)). Substituting
Theorem 7.2 into (E\_4):

{[}
4\int\emph{\{N+1/2\}\textsuperscript{\infty \frac{\Delta^*(x)}{x}\Xi\emph{B(T,x),dx
=
-\frac{8}{\pi}\sum}\{m=1\}}\infty \frac{d(m)}{\sqrt{m}}\int}\{N+1/2\}\textsuperscript{\infty \frac{K_1(4\pi\sqrt{mx}) + \frac{\pi}{2}Y_1(4\pi\sqrt{mx})}{\sqrt{x}}\Xi\emph{B(T,x),dx
{]} {[}
+\frac{16\sqrt{2}}{\pi}\sum}\{m=1\}}\infty \frac{d(m)}{\sqrt{m}}\int\emph{\{N+1/2\}\textsuperscript{\infty \frac{K_1(4\pi\sqrt{2mx}) + \frac{\pi}{2}Y_1(4\pi\sqrt{2mx})}{\sqrt{x}}\Xi\emph{B(T,x),dx
{]} {[}
-\frac{16}{\pi}\sum}\{m=1\}}\infty \frac{d(m)}{\sqrt{m}}\int}\{N+1/2\}\^{}\infty \frac{K_1(4\pi\sqrt{4mx}) + \frac{\pi}{2}Y_1(4\pi\sqrt{4mx})}{\sqrt{x}}\Xi\_B(T,x),dx.
\tag{6.8} {]}

Substituting (6.8) into (6.1) gives a complete representation of (E(T))
as a convergent series in (K\_0, Y\_0, K\_1, Y\_1) with no error term
other than (R\_0(T)) from (5.4).

\emph{Proof.} Substitute (7.2) into (5.8) term by term. The three lines
of (6.8) correspond to the three terms (-\Delta(x)), (+2\Delta(2x)),
(-\tfrac{1}{2}\Delta(4x)) in (\Delta\^{}*), with the Voronoï series
(7.1) substituted for each. The prefactors are:
(-(-\tfrac{2}{\pi})\cdot\tfrac{4}{\sqrt{x}} = \tfrac{8}{\pi\sqrt{x}})
for the first,
(+2\cdot(-\tfrac{2}{\pi})\cdot(-\tfrac{4\sqrt{2}}{\sqrt{x}}) =
\tfrac{16\sqrt{2}}{\pi\sqrt{x}}) for the second,
(-(-\tfrac{1}{2})\cdot(-\tfrac{2}{\pi})\cdot\tfrac{8}{\sqrt{x}} =
-\tfrac{16}{\pi\sqrt{x}}) for the third. \(\square\)

\begin{center}\rule{0.5\linewidth}{0.5pt}\end{center}

\subsection{\texorpdfstring{7. Voronoï's Exact Series for
(\Delta\^{}*(x))}{7. Voronoï's Exact Series for (\^{}*(x))}}\label{voronouxefs-exact-series-for-x}

\textbf{Theorem 7.1} (Voronoï 1904 --- exact series for (\Delta(x))).
For (x \textgreater{} 0), (x \notin \mathbb{Z}):

{[} \Delta(x) =
-\frac{2}{\pi}\sqrt{x}\sum\_\{n=1\}\^{}\infty \frac{d(n)}{\sqrt{n}}\left[K_1(4\pi\sqrt{nx}) + \frac{\pi}{2}Y_1(4\pi\sqrt{nx})\right].
\tag{7.1} {]}

The series converges absolutely in (K\_1) and conditionally in (Y\_1),
uniformly on compact subsets of ((0,\infty)\setminus\mathbb{Z}). This is
an exact identity.

\emph{Proof.} Apply the Mellin transform to (\Delta(x)), shift the
inverse Mellin contour to (\mathrm{Re}(s) = \tfrac{1}{2}) via the
functional equation of (\zeta(s)\^{}2), and identify the resulting
integral with the Bessel functions via the Basset formulas:

{[} K\_1(2\pi a) =
\int\_0\textsuperscript{\infty e}\{-2\pi a\cosh t\}\cosh t,dt,
\qquad Y\_1(2\pi a) =
-\frac{2}{\pi}\int\_0\^{}\infty \cosh t\cdot\sin(2\pi a\cosh t),dt. {]}

Summing over (n) with coefficients (d(n)/\sqrt{n}) gives (7.1).
\(\square\)

\textbf{Theorem 7.2} (Exact series for (\Delta\^{}\emph{(x))). Let
(\Delta\^{}}(x) = -\Delta(x) + 2\Delta(2x) - \tfrac{1}{2}\Delta(4x)).
Substituting (7.1) into each term:

{[} \Delta\^{}*(x) =
\frac{2}{\pi}\sqrt{x}\sum\emph{\{n=1\}\^{}\infty \frac{d(n)}{\sqrt{n}}\left[K_1(4\pi\sqrt{nx}) + \frac{\pi}{2}Y_1(4\pi\sqrt{nx})\right]{]}
{[} -
\frac{4\sqrt{2}}{\pi}\sqrt{x}\sum}\{n=1\}\^{}\infty \frac{d(n)}{\sqrt{n}}\left[K_1(4\pi\sqrt{2nx}) + \frac{\pi}{2}Y_1(4\pi\sqrt{2nx})\right]{]}
{[} +
\frac{2}{\pi}\sqrt{x}\sum\_\{n=1\}\^{}\infty \frac{d(n)}{\sqrt{n}}\left[K_1(4\pi\sqrt{4nx}) + \frac{\pi}{2}Y_1(4\pi\sqrt{4nx})\right].
\tag{7.2} {]}

\emph{Proof.} Line-by-line substitution of (7.1): (-\Delta(x)) gives
coefficient (+(2/\pi)\sqrt{x}); (2\Delta(2x)) gives
(+2\cdot(-2/\pi)\sqrt{2x} = -(4\sqrt{2}/\pi)\sqrt{x});
(-\tfrac{1}{2}\Delta(4x)) gives
(-\tfrac{1}{2}\cdot(-2/\pi)\cdot 2\sqrt{x} = +(2/\pi)\sqrt{x}).
\(\square\)

\textbf{Corollary 7.3} (Jutila's mean-square relation). Jutila proved:

{[} \int\_2\^{}T
\left[E(t) - 2\pi\Delta^*\!\left(\frac{t}{2\pi}\right)\right]\^{}2 dt =
O(T\textsuperscript{\{4/3\}\log}3 T). {]}

This shows (E(T) \approx 2\pi\Delta\^{}*(T/(2\pi))) in the (L\^{}2)
sense. Substituting (7.2) at (x = T/(2\pi)) gives the leading Bessel
series for (E(T)); the exact pointwise formula is (6.1).

\begin{center}\rule{0.5\linewidth}{0.5pt}\end{center}

\subsection{8. Convergence}\label{convergence}

\textbf{Theorem 8.1.}

\textbf{(i)} The finite sum (\sum\_\{n=1\}\^{}N d(n)\Psi\_B(T,n))
contains (N \asymp T) terms, each an explicit finite integral. It
converges trivially for fixed (T).

\textbf{(ii)} (\Xi\_B(T,x) = O(Tx\^{}\{-1/2\})) as (x \to \infty) by the
Riemann--Lebesgue lemma. Combined with (\Delta\^{}*(x) =
O(x\^{}\{1/3+\varepsilon\})):

{[}
\int\_\{N+1/2\}\^{}\infty \frac{|\Delta^*(x)|}{x}\textbar{}\Xi\_B(T,x)\textbar,dx
\le CT N\^{}\{-1/6+\varepsilon\} \textless{} \infty. {]}

\textbf{(iii)} In the double Bessel series (6.8): the (K\_1) terms
satisfy (\textbar d(m)K\_1(4\pi\sqrt{mx})\textbar{}
\le Cm\textsuperscript{\varepsilon e}\{-4\pi\sqrt{mx}\}), giving
absolute and exponentially fast convergence. The (Y\_1) terms converge
conditionally by Dirichlet's criterion after Abel summation on (d(m)).

\textbf{(iv)} For any (T \ge 2) and target precision
(\delta \textgreater{} 0), truncating the Voronoï series in (6.8) at (m
= M\_0) gives

{[} \textbar E(T) - E(T;M\_0)\textbar{} \textless{} \delta, {]}

with the (K\_1) tail bounded by (C',e\^{}\{-4\pi\sqrt{M_0(N+1/2)}\}),
which is effective and computable.

\begin{center}\rule{0.5\linewidth}{0.5pt}\end{center}

\subsection{\texorpdfstring{9. The Saddle-Point Step and the Origin of
(O(\log\^{}2
T))}{9. The Saddle-Point Step and the Origin of (O(\^{}2 T))}}\label{the-saddle-point-step-and-the-origin-of-o2-t}

\textbf{Proposition 9.1} (Stationary-phase lemma). Let (\varphi, g :
{[}a,b{]} \to \mathbb{R}) with (\varphi) smooth, having a unique
stationary point (x\_0 \in (a,b)) with (\varphi`(x\_0) = 0),
(\varphi'\,'(x\_0) \ne 0), and (g = O(G)). Then:

{[} \int\_a\^{}b g(x)e\^{}\{i\lambda\varphi(x)\},dx =
g(x\_0)e\textsuperscript{\{i\lambda\varphi(x\_0)\}\sqrt{\frac{2\pi}{\lambda|\varphi''(x_0)|}},e}\{i(\pi/4)\mathrm{sgn},\varphi'\,'(x\_0)\}
+ O!\left(\frac{G}{\lambda|\varphi''(x_0)|}\right). \tag{9.1} {]}

\textbf{Proposition 9.2.} The integrand of (\Psi\_B(T,n)) has phase
(\varphi(y) = T\log((1+y)/y) \pm 2\pi ny). Setting
(\varphi'(y\textsuperscript{\emph{) = 0) gives (T/(y\^{}}(1+y}\emph{)) =
2\pi n), with unique solution (y\^{}} \approx \sqrt{T/(4\pi n)} -
\tfrac{1}{2}) for (T \gg n). Applying (9.1):

{[} \Psi\_B(T,n) =
\frac{\pi}{2}\left(\frac{T}{2\pi n}\right)\textsuperscript{\{1/4\}!\left(\frac{T}{2\pi n}-1\right)}\{-1/2\}\cos(f(T,n))
+ O(n\textsuperscript{\varepsilon T}\{-3/4\}), {]}

where

{[} f(T,n) = T\operatorname{arcsinh}!\sqrt{\frac{2\pi n}{T}} -
\sqrt{2\pi nT + 4\pi^2 n^2} + \frac{\pi}{4}. \tag{9.2} {]}

Summing over (n = 1,\ldots,N) with weight (d(n)): the total saddle-point
error is (\sum\_\{n=1\}\^{}N
d(n)O(n\textsuperscript{\varepsilon T}\{-3/4\}) =
O(T\^{}\{1/4+2\varepsilon\})). The boundary terms (E\_2) and (E\_3),
evaluated at the large-argument boundary (n = N+\tfrac{1}{2} \asymp T),
contribute (O(\log\^{}2 T)) from the stationary-phase remainder at that
scale. This is the entire source of the published (O(\log\^{}2 T)) in
the standard Atkinson formula. Without applying (9.1), formula (6.1)
carries no error beyond (R\_0(T)).

\textbf{Corollary 9.3} (Standard Atkinson formula --- approximate).
Applying Proposition 9.1 to each (\Psi\_B(T,n)) and (\Xi\_B(T,x)):

{[} E(T) = \sum\emph{\{n \le N\} a(n)\cos(f(T,n)) + \sum}\{n \le M\}
b(n)\cos(g(T,n)) + O(\log\^{}2 T), \tag{9.3} {]}

with

{[} a(n) = (-1)\^{}n
d(n)\left(\frac{T}{2\pi n}\right)\^{}\{1/4\}!\left(1 -
\frac{2\pi n}{T}\right)\^{}\{-1/4\}, \qquad b(n) =
\frac{1}{\sqrt{2}}d(n)\left(\frac{2T}{\pi n}\right)\^{}\{1/4\}!\left(1 +
\frac{\pi n}{2T}\right)\^{}\{-1/4\}, {]}

{[} g(T,n) = -2T\operatorname{arcsinh}!\sqrt{\frac{\pi n}{2T}} +
\sqrt{2\pi nT + \pi^2 n^2} + \frac{\pi}{4}. {]}

The (O(\log\^{}2 T)) is entirely a consequence of applying (9.1). It
does not appear in (6.1).

\begin{center}\rule{0.5\linewidth}{0.5pt}\end{center}

\subsection{10. Summary of Identities}\label{summary-of-identities}

\begin{longtable}[]{@{}
  >{\raggedright\arraybackslash}p{(\linewidth - 4\tabcolsep) * \real{0.2800}}
  >{\raggedright\arraybackslash}p{(\linewidth - 4\tabcolsep) * \real{0.4000}}
  >{\raggedright\arraybackslash}p{(\linewidth - 4\tabcolsep) * \real{0.3200}}@{}}
\toprule\noalign{}
\begin{minipage}[b]{\linewidth}\raggedright
Label
\end{minipage} & \begin{minipage}[b]{\linewidth}\raggedright
Identity
\end{minipage} & \begin{minipage}[b]{\linewidth}\raggedright
Status
\end{minipage} \\
\midrule\noalign{}
\endhead
\bottomrule\noalign{}
\endlastfoot
((*)) & (\int\_0\^{}T & \zeta(\tfrac{1}{2}+it) \\
(3.1) & Voronoï--Jutila summation formula & \textbf{Exact} \\
(4.1) & Titchmarsh identity for (\zeta\^{}2(s)) with (K\_1), (Y\_1)
Bessel tail & \textbf{Exact} \\
(5.9) & Atkinson decomposition (E(T) = E\_1 - E\_2 + E\_3 - E\_4 +
R\_0(T)) & \textbf{Exact} \\
(6.1) & Main theorem: (E(T)) as explicit Bessel integrals &
\textbf{Exact} \\
(6.8) & (E\_4) as double series in (K\_1, Y\_1) & \textbf{Exact} \\
(7.1) & Voronoï exact series for (\Delta(x)) & \textbf{Exact} \\
(7.2) & Exact series for (\Delta\^{}*(x)) & \textbf{Exact} \\
(9.3) & Standard Atkinson formula & \textbf{Approximate ((O(\log\^{}2
T)) error)} \\
\end{longtable}

\end{document}
